\chapter{Macrofinanzas con agentes heterogéneos}
\label{cap:macrofinanzas}

Casi toda la literatura de jacobianos secuenciales está escrita con hogares como
agentes heterogéneos. Pero nada del método lo exige: el ``agente'' puede ser una
empresa, un intermediario o un banco. Este capítulo recorre lo que cambia
cuando el agente es una institución financiera, porque hay varias diferencias
sustantivas y no todas son obvias.

\section{Qué cambia cuando el agente es un banco}

\paragraph{El estado endógeno es el patrimonio, no la riqueza} Y su ley de
movimiento es distinta: el patrimonio de un banco crece con beneficios retenidos
y cae con dividendos y pérdidas. La restricción relevante no suele ser ``no
puede endeudarse'' sino ``no puede apalancarse más allá de un múltiplo de su
patrimonio''.

\paragraph{Los retornos son aproximadamente lineales en el patrimonio} Un banco
que duplica su capital duplica sus depósitos, sus créditos y sus beneficios. Esa
homoteticidad tiene una consecuencia incómoda: si el modelo es exactamente
homotético, la distribución de tamaños no afecta a los agregados, y el bloque
heterogéneo colapsa a un conjunto de bancos representativos, uno por tipo.

\begin{idea}
Para que la distribución de tamaños \emph{haga trabajo} en un modelo bancario
hace falta alguna no homoteticidad. Las tres habituales son: un costo operativo
fijo, una restricción de apalancamiento que dependa del tamaño, o un acceso al
mercado mayorista que dependa del tamaño. La tercera es la más defendible
empíricamente --- los bancos pequeños tienen peor acceso al interbancario y al
mercado de capitales --- y es la que usamos en la parte V.
\end{idea}

\paragraph{La entrada y la salida son de primer orden} En modelos de hogares, la
entrada y salida suelen ser un detalle demográfico. En modelos bancarios son el
mecanismo que fija la escala: sin salida, un banco con retornos superiores a la
tasa de descuento crece sin límite. El dispositivo estándar --- los banqueros
sobreviven con probabilidad $\varsigma$ y son reemplazados por entrantes de
tamaño fijo --- viene de Gertler y Kiyotaki (2010) y Gertler y Karadi (2011), y
encaja exactamente en la proposición \ref{prop:general}.

\paragraph{Las fricciones de pago de dividendos sustituyen a la aversión al
riesgo} Un banquero neutral al riesgo con retornos lineales pagaría todo o nada
como dividendo. Para obtener una política interior hace falta alguna
concavidad: utilidad cóncava sobre dividendos, costos de emisión de capital, o
un valor de franquicia. Todas son formas reducidas de la misma fricción
financiera, y conviene ser explícito sobre cuál se usa y qué se pierde.

\section{Riesgo agregado dentro del período}

Este es el punto técnico más importante del capítulo, y vale para cualquier
modelo macrofinanciero.

El método es de primer orden en los agregados: por equivalencia cierta, el riesgo
agregado desaparece. Pero en macrofinanzas el riesgo agregado \emph{es} el
objeto de estudio: una crisis de liquidez, una salida de capitales, un evento de
cola. ¿Cómo se concilian?

La respuesta es distinguir dos tipos de incertidumbre agregada.

\begin{enumerate}[leftmargin=1.4em]
\item \textbf{Incertidumbre entre períodos.} ``No sé si el año que viene habrá
una crisis''. Esta desaparece a primer orden. No hay forma de recuperarla sin ir
a orden superior o a métodos globales.
\item \textbf{Incertidumbre dentro del período.} ``Sé que en cada período, con
probabilidad $p$, ocurre un episodio de escasez de dólares, y conozco sus
consecuencias''. Esta \textbf{sí} se puede conservar de forma completamente no
lineal.
\end{enumerate}

La forma de conservar la segunda es tratar el estado agregado intraperiodo
$A\in\{A_1,\dots,A_m\}$ como una dimensión adicional sobre la que se integra
\emph{dentro} del paso hacia atrás, con probabilidades fijas. Las variables
agregadas que dependen del estado --- un precio de mercado que se vacía, una
capacidad de respaldo que se raciona --- se convierten en $m$ secuencias
separadas, cada una con su propia condición de equilibrio.

\begin{receta}[Incorporar un estado agregado intraperiodo]
\begin{enumerate}[leftmargin=1.4em,itemsep=2pt]
\item Define el conjunto finito de estados $A$ y sus probabilidades, que
permanecen fijas (o son choques exógenos del modelo).
\item Para cada variable agregada que se determine \emph{condicional} en $A$
--- tasas de mercados que se vacían, fracciones racionadas, coberturas ---
declara una secuencia por estado: $\Psi_t(A_1),\dots,\Psi_t(A_m)$.
\item En el paso hacia atrás, el agente toma su decisión \emph{ex ante},
integrando sobre $A$ con las probabilidades dadas. La condición de primer orden
pondera los estados.
\item En el DAG, cada $\Psi(A_j)$ es una incógnita y su condición de vaciado en
el estado $j$ es un objetivo.
\end{enumerate}
\end{receta}

El costo es que el número de incógnitas se multiplica por $m$. El beneficio es
que el modelo conserva exactamente la no linealidad del evento de cola: la
diferencia entre un estado normal y uno de crisis no está linealizada.

\section{Quiebres, racionamiento y esquinas}

Los modelos macrofinancieros están llenos de $\min$ y de $\max$: capacidad de
respaldo limitada, racionamiento proporcional, restricciones de colateral que
atan o no atan. La sección \ref{sec:quiebres} dio la regla general. Tres
consecuencias prácticas:

\begin{itemize}[leftmargin=1.4em]
\item \textbf{Calibra lejos del quiebre, y reporta cuán lejos.} Si el racionamiento
en el estado estacionario es $P/Q = 0.36$, hay mucho espacio antes de llegar a
$1$. Si fuera $0.97$, el modelo lineal sería engañoso.
\item \textbf{Las esquinas individuales son distintas de los quiebres
agregados.} Que un banco esté en la esquina $u=0$ no es un problema: la
derivada existe y vale cero localmente. Que el agregado esté en un quiebre sí lo
es.
\item \textbf{Usa el capítulo \ref{cap:nolineal} para medir el daño.} Si el
quiebre está a dos desviaciones estándar, conviene calcular la transición no
lineal a un choque de ese tamaño y mostrar cuánto se separa.
\end{itemize}

\section{Qué preguntas se vuelven contestables al pasar a dinámico}

Vale la pena cerrar con esto, porque es la justificación del esfuerzo. Un modelo
estático con heterogeneidad puede responder preguntas de \emph{incidencia en un
punto}: quién usa más la cobertura, quién paga más del costo. Al pasar al espacio
de secuencias se abren cuatro familias de preguntas nuevas:

\begin{enumerate}[leftmargin=1.4em]
\item \textbf{Transición versus estado estacionario.} Una reforma puede ser
beneficiosa en el largo plazo y costosa durante la transición, y la transición
puede durar décadas. El modelo estático no distingue.
\item \textbf{Anticipación.} ¿Qué pasa cuando la política se \emph{anuncia} con
antelación? Las columnas del jacobiano responden esto directamente, y en
regulación financiera --- donde los períodos de implementación gradual son la
norma --- es una pregunta de política de primer orden.
\item \textbf{Acumulación y erosión de capital.} El costo de una política que
reduce la rentabilidad bancaria no es solo el efecto de impacto: es la erosión
acumulada del patrimonio, que reduce el crédito por años.
\item \textbf{Retroalimentación entre heterogeneidad y agregados.} La
composición del sistema cambia con el tiempo en respuesta a la política. Un
encaje que castiga desproporcionadamente a cierto tipo de banco cambia la
participación de ese tipo, y eso a su vez cambia la respuesta agregada futura.
Esto es intrínsecamente dinámico.
\end{enumerate}
